\section{Hai thuật toán gia tốc}

\subsection{Cấu trúc chung}

Dãy ước lượng bậc hai được viết dưới dạng
\[
\varphi_k(x)=(\varphi_k)_{\star}
+\frac{A_k}{2}\|x-\nu_k\|^2,
\qquad
A_{k+1}=(1-\alpha_k)A_k.
\]
Với $\beta_0=1$, đặt
\[
\beta_k=\prod_{i=0}^{k-1}(1-\alpha_i).
\]
Các quy tắc chọn $\alpha_k$ ở hai thuật toán bảo đảm tốc độ nền $\beta_k=O(k^{-2})$ dưới giả thiết so sánh của dãy tham số $(\lambda_k)$.

\subsection{IAPPA1 với xấp xỉ loại 1}

Tiêu chuẩn loại 1 không cung cấp trực tiếp phần tử dưới vi phân dùng để cập nhật tâm của dãy ước lượng. IAPPA1 thay tâm lý tưởng $\nu_k$ bằng một biến tính được $u_k$, đồng thời dùng $\eta_k$ để chặn sai lệch giữa hai đại lượng.

\medskip
\noindent\textbf{Thuật toán IAPPA1.}
Khởi tạo $x_0\in\operatorname{dom}F$, $u_0=x_0$, $A_0=A>0$ và $\eta_0=\delta_0=0$. Với mỗi $k\geq0$:
\begin{align*}
\alpha_k^2&=(1-\alpha_k)A_k\lambda_k,\\
y_k&=(1-\alpha_k)x_k+\alpha_k u_k,\\
x_{k+1}&\approx_1\prox{\lambda_kF}(y_k)
\quad\text{với độ chính xác }\varepsilon_k,\\
A_{k+1}&=(1-\alpha_k)A_k,\\
u_{k+1}&=u_k-\frac{1}{\alpha_k}(y_k-x_{k+1}),\\
\eta_{k+1}&=\eta_k+\frac{\varepsilon_k}{\alpha_k},\\
\delta_{k+1}&=(1-\alpha_k)\delta_k
+\frac{(\alpha_k\eta_{k+1})^2}{2\lambda_k}.
\end{align*}
Các đại lượng sai số có dạng tường minh
\begin{equation}
\eta_k=\sum_{i=0}^{k-1}\frac{\varepsilon_i}{\alpha_i},
\qquad
\delta_k=\frac{A\beta_k}{2}
\sum_{i=0}^{k-1}\eta_{i+1}^2.
\label{eq:iappa1-error}
\end{equation}
Mỗi sai số $\varepsilon_i$ được khuếch đại bởi $1/\alpha_i$, cộng dồn trong $\eta_k$, rồi các tổng riêng này mới được bình phương trong $\delta_k$.

\subsection{IAPPA2 với xấp xỉ loại 2}

Tiêu chuẩn loại 2 cung cấp trực tiếp phần tử dưới vi phân
\[
\frac{y_k-x_{k+1}}{\lambda_k}
\in\partial_{\varepsilon_k^2/(2\lambda_k)}F(x_{k+1}).
\]
Vì vậy, IAPPA2 cập nhật tâm $\nu_k$ mà không cần biến sai lệch $\eta_k$.

\medskip
\noindent\textbf{Thuật toán IAPPA2.}
Khởi tạo $x_0=\nu_0\in\operatorname{dom}F$, $A_0=A>0$, $\delta_0=0$ và chọn $a\in(0,2]$. Với mỗi $k\geq0$:
\begin{align*}
a&\leq\frac{\alpha_k^2}{(1-\alpha_k)A_k\lambda_k}\leq2,\\
y_k&=(1-\alpha_k)x_k+\alpha_k\nu_k,\\
x_{k+1}&\approx_2\prox{\lambda_kF}(y_k)
\quad\text{với độ chính xác }\varepsilon_k,\\
A_{k+1}&=(1-\alpha_k)A_k,\\
\nu_{k+1}&=\nu_k-
\frac{\alpha_k}{(1-\alpha_k)A_k\lambda_k}(y_k-x_{k+1}),\\
\delta_{k+1}&=(1-\alpha_k)\delta_k+
\frac{\varepsilon_k^2}{2\lambda_k}.
\end{align*}
Nghiệm của truy hồi sai số là
\begin{equation}
\delta_k=\frac{\beta_k}{2}
\sum_{i=0}^{k-1}
\frac{\varepsilon_i^2}{\lambda_i\beta_{i+1}}.
\label{eq:iappa2-error}
\end{equation}
Công thức \eqref{eq:iappa2-error} chỉ chứa các bình phương riêng lẻ của sai số. Nó không chứa bình phương của một tổng sai số đã tích lũy như \eqref{eq:iappa1-error}.

